\documentclass[10pt,a4paper]{amsart}
\usepackage[margin=19mm]{geometry}
\usepackage{amsmath,amssymb,mathtools}
\usepackage[scr=rsfs,cal=euler]{mathalfa}
\usepackage{xfrac}
\usepackage[unicode,hidelinks,bookmarks=false]{hyperref}
\newcommand{\ie}{i.e.\ }
\newcommand{\cf}{cf.\ }
\newcommand{\resp}{respectively\ }
\newcommand{\Subsection}{\subsection}
\providecommand{\nobreakdash}{\nobreak\hbox{-}}
\setlength{\emergencystretch}{2em}
\multlinegap0pt
\usepackage{bm}
\usepackage{bbm}
\usepackage{dsfont}
\usepackage{xspace}
\usepackage{cancel}
\usepackage{mathtools}
\usepackage{amsthm}
\makeatletter
\@ifundefined{thm}{\newtheorem{thm}{Thm}}{}
\@ifundefined{defn}{\newtheorem{defn}[thm]{Definition}}{}
\@ifundefined{lemma}{\newtheorem{lemma}[thm]{Lemma}}{}
\@ifundefined{prop}{\newtheorem{prop}[thm]{Proposition}}{}
\makeatother
\newcommand{\Aut}{\operatorname{Aut} }
\newcommand{\IM}{\operatorname{im} }
\newcommand{\Ker}{\operatorname{ker} }
\title[Cohomology and Extensions of Relative Rota-Baxter groups]{Cohomology and Extensions of Relative Rota-Baxter groups}
\author{Pragya Belwal, Nishant Rathee, Mahender Singh}
\date{Source: arXiv:2309.00692v1 (CC BY 4.0)}
\begin{document}
\maketitle

\noindent\emph{Source and licence.} Cohomology and Extensions of Relative Rota-Baxter groups. Version arXiv:2309.00692v1 (CC BY 4.0), distributed under the Creative Commons Attribution 4.0 International licence, as recorded in the arXiv licence metadata for this exact version and shown on the abstract page https://arxiv.org/abs/2309.00692v1. Full rights notice: licenses/arxiv-2309.00692-CC-BY-4.0.txt.\par\smallskip

\noindent\textit{Editorial scope.} Counted unit: the Lemma whose source label is \texttt{rrb coboundary condition} in the article (source0.tex lines 760--762, source label \texttt{rrb coboundary condition}), delivered together with the article's own complete original proof of it (source0.tex lines 764--813, one \texttt{proof} environment of 50 lines). No mathematics is written, repaired or re-derived: statement and proof are the article's own text line for line. Required context retained uncounted: group coboundary mu (lines 284--287); group coboundary twisted sigma (lines 292--295); phidefnprop (lines 567--569); Rallcondn (lines 652--665); moddefn (lines 688--697). Every definition, display and label of those lines is retained unchanged. Omitted from this package and not redistributed: 1--283, 288--291, 296--566, 570--651, 666--687, 698--759, 763--763, 814--1311. Nothing inside a retained excerpt is omitted or altered: every retained body is delivered exactly as the article's lines are written, so no omission receipt of any kind accompanies this package. Citations: the retained excerpts carry no citation command at all, so no bibliography entry of the article is transcribed and no reference is redistributed. Reference adaptation: none was needed and none was made; every reference inside the delivered unit resolves inside it, so no unresolved reference or citation is produced. Printed slips kept literal: no printed word, symbol, hypothesis or number of the counted statement or of its proof was altered. Layout and class: the article's own class is \texttt{amsart} with packages including amsmath, amsfonts, amssymb, upgreek, amscd, mathtools, hyperref, tikz, enumerate, tikz-cd, graphicx; the wrapper is the lane's pinned \texttt{amsart} editorial wrapper at 10pt on A4, which restates the theorem-like environments the retained text needs and adds the article's own macro definitions that the retained text needs (\texttt{\textbackslash Aut}, \texttt{\textbackslash IM}, \texttt{\textbackslash Ker}), with extra wrapper packages (bm, bbm, dsfont, xspace, cancel, mathtools). The delivered numbering is the unit's own. Assets: no retained span contains a figure environment, a graphics-inclusion command, a PDF-page inclusion, a table or a TikZ picture; no image file is copied into this package, the package declares no asset dependency and no image file is redistributed with it. Licence: version arXiv:2309.00692v1 of this article is distributed under the Creative Commons Attribution 4.0 International licence, as recorded in the arXiv licence metadata for this exact version and shown on the abstract page https://arxiv.org/abs/2309.00692v1. A full rights notice is delivered at licenses/arxiv-2309.00692-CC-BY-4.0.txt. Characters: every retained excerpt is pure ASCII, so the delivered engine, XeLaTeX with its default Latin Modern text font, reports no missing character.
\begin{eqnarray}\label{group coboundary mu}
&& (\partial_\mu^n f)(a_1, a_2,\ldots, a_{n+1}) \\
&= & f(a_2, \ldots,a_{n+1})~ \prod^{n}_{k=1}    (f(a_1, \ldots, a_k \cdot a_{k+1}, \ldots, a_{n+1}))^{(-1)^k}~ \mu_{a_{n+1}}(f(a_1, \ldots,a_{n})), \notag
\end{eqnarray} 

\begin{eqnarray}\label{group coboundary twisted sigma}
&&	(\delta_{\sigma}^n f)(a_1, a_2,\ldots, a_{n+1})\\
 &= & f(a_2, \ldots,a_{n+1}) ~\prod^{n}_{k=1}    (f(a_1, \ldots, a_k \circ_R a_{k+1}, \ldots, a_{n+1}))^{(-1)^k} ~ \sigma_{R(a_{n+1})}(f(a_1, \ldots,a_{n})). \notag
\end{eqnarray}

		\begin{eqnarray}\label{phidefnprop}
			\phi_{(b,l)}(a,k) &= & \big(\beta{_{b}(a)}, ~\rho(a,b)\nu_b(f(l,a)k)\big),
		\end{eqnarray} 

\begin{prop}\label{Rallcondn}
Let $\phi: B \times_{\tau_2} L \longrightarrow \Aut(A \times_{\tau_1} K )$ defined in \eqref{phidefnprop} be a homomorphism. Then the map  $R:  A \times_{\tau_1} K \rightarrow B \times_{\tau_2} L$ given by
\begin{eqnarray*}
R(a,k) &= & \big(T(a), ~\chi(a) \,S(\nu^{-1}_{T(a)}(k))\big),
\end{eqnarray*} 
for $(a,k) \in A \times_{\tau_1} K$ and $(b,l) \in B \times_{\tau_2} L$, is a relative Rota-Baxter operator if and only if the conditions 
\begin{eqnarray*}
& & \chi(a_1 \circ_{T} a_2) \,S \big(\nu^{-1}_{T(a_1 \circ_{T} a_2)}\big(\rho(a_2, T(a_1)) \,\tau_1(a_1,\beta{_{T(a_1)}}(a_2))\big) \,\nu^{-1}_{T(a_2)}(f(\chi(a_1), a_2))\big) \\
& =&   \sigma_{T(a_2)}( \chi(a_1)) \, \chi(a_2) \,\tau_2(T(a_1), T(a_2) )\\
\textrm{and} &&\\
&& S\big(\nu^{-1}_{T(a_2)}(\mu_{a_2}(k)) \,\nu^{-1}_{T( a_2)}(f(S(k), a_2))\big) = \sigma_{T(a_2)}(S(k)),	
\end{eqnarray*}		
hold for all $ a_1, a_2 \in A$ and $ k \in K$.
\end{prop}

\begin{defn}\label{moddefn}
	A module over a relative Rota-Baxter group $(A, B,\beta, T)$ is a trivial relative Rota-Baxter group $(K, L, \alpha, S)$ such that there exists a quadruple $(\nu, \mu, \sigma, f)$ (called action) of maps satisfying the following conditions:
	\begin{enumerate}
		\item The group $K$ is a left $B$-module and a right $A$-module with respect to the actions $\nu:B \to \Aut(K) $ and $\mu: A \to \Aut(K)$, respectively.
		\item The group $L$ is a right $B$-module with respect to the action $\sigma: B \to \Aut(L)$.
		\item The  map  $f: L \times A \to K$ has the property that $f(-,a): L \rightarrow K$ is a homomorphism for all $a \in A$ and $f(l,-): A \rightarrow K$ is a derivation with respect to the action $\mu$ for all $l \in L$.
		\item $S\big(\nu^{-1}_{T(a)}(\mu_{a}(k)) \,\nu^{-1}_{T( a)}(f(S(k), a))\big)=\;\sigma_{T(a)}(S(k))$ for all $ a \in A$ and $k \in K$.
		\item $\nu_b(\mu_{a}(k))= \;   \mu_{ \beta{_{b}(a)}}( \nu_{b} (k))$ for all $a \in A$, $b \in B$ and $k \in K$. 
	\end{enumerate}
\end{defn}

\begin{lemma}\label{rrb coboundary condition}
	$\IM(\partial^0_{RRB}) \subseteq  \Ker(\partial^1_{RRB})$ and	$\IM(\partial^1_{RRB}) \subseteq  \Ker(\partial^2_{RRB})$.
\end{lemma}

\begin{proof}
It is enough to prove that $\partial^1_{RRB}\;  \partial^0_{RRB}$ and $\partial^2_{RRB}\;  \partial^1_{RRB}$ are trivial homomorphisms. We have 
	\begin{eqnarray*}
		\partial^1_{RRB}\;  \partial^0_{RRB}(k,l) &=& (\partial^1_{\mu}(\kappa_k), \partial^1_{\sigma}(\omega_l), \lambda_1, \lambda_2).
	\end{eqnarray*}
	Since $\kappa_k$ and $\omega_l$ are group 1-coboundaries,  we have $\partial^1_{\mu}(\kappa_k)=1$ and $\partial^1_{\sigma}(\omega_l)=1$. Further, we have
	\begin{eqnarray*}
		\lambda_1(a,b) &=& \nu_b \big(f(\sigma_b(l)l^{-1}, a)\mu_a(k)k^{-1} \big)  \,(\mu_{\beta_b(a)}(k))^{-1} \,k.
	\end{eqnarray*}
	Using the definition of $K \times_{(\nu, \mu, \sigma, f)} L$ and the compatibility conditions in Definition \ref{moddefn}, we have $\lambda_1(a, b) = 1$. Similarly, we obtain $\lambda_2(a, b) = 1$ for all $a \in A$ and $b \in B.$ 
		\par 
	
	Now we prove that $\partial^2_{RRB}\;  \partial^1_{RRB}$ is trivial. If $(\theta_1, \theta_2) \in C^1_{RRB}$, then
	\begin{eqnarray}\label{com1}
		\partial^2_{RRB}\;  \partial^1_{RRB}(\theta_1, \theta_2) &=& \partial^2_{RRB}(\partial_{\mu}^1(\theta_1), \partial_{\sigma}^1(\theta_2), \lambda_1, \lambda_2 ).
	\end{eqnarray}
	We show that each term on the right hand side of \eqref{com1} is a trivial map. The first two terms are trivial maps since $\partial_{\mu}^2\partial_{\mu}^1$ and $  \partial_{\sigma}^2\partial_{\sigma}^1$ are trivial being the coboundary maps defining the usual group cohomology. We expand the remaining three terms using the definitions of $\partial^1_{RRB}$ and $\partial^2_{RRB}$. It follows from \eqref{group coboundary mu} and \eqref{group coboundary twisted sigma} that
	\begin{eqnarray*}
		\partial^1_{\mu}(\theta_1)(a_1, a_2) &=& \theta_1(a_2)\, (\theta_1(a_1 a_2))^{-1} \,\mu_{a_2}(\theta_1(a_1)),\\
		\partial^1_{\sigma}(\theta_2)(b_1, b_2) &=& \theta_2(b_2) \,(\theta_2(b_1 b_2))^{-1} \,\sigma_{b_2}(\theta_2(b_1)),\\
		\delta^1_{\sigma}(\chi)(a_1, a_2) &=& \chi(a_2) \, (\chi(a_1 \circ_T a_2))^{-1} \,\sigma_{T(a_2)}( \chi(a_1)).
	\end{eqnarray*}
	
	The third term of \eqref{com1} is given by
	\begin{eqnarray}\label{t1}
		&& \gamma_1(a, b_1, b_2) \notag\\
		&=&	  \nu_{b_1 b_2}(f(\theta_2(b_1 b_2), a)\theta_1(a))(\theta_1(\beta_{b_1 b_2}(a)))^{-1} \nu_{b_1 b_2}(f( \theta_2(b_2)(\theta_2(b_1 b_2))^{-1}\sigma_{b_2}(\theta_2(b_1)), a))\notag \\ 
		&&(\nu_{b_1}(f(\theta_2(b_1), \beta_{b_2}(a))\theta_1(\beta_{b_2}(a)))(\theta_1(\beta_{b_1}(\beta_{b_2}(a))))^{-1})^{-1}(\nu_{b_1}\big( \nu_{b_2}(f(\theta_2(b_2), a)\theta_1(a))(\theta_1(\beta_{b_2}(a)))^{-1})\big)^{-1}. \notag
	\end{eqnarray}
Using that $f(\theta_2(b_1), \beta_{b_2}(a))= \nu_{b_2}(f(\sigma_{b_2}(\theta_2(b_2)),a))$ and $f$ is linear in the first coordinate, we have $\gamma_1(a, b_1, b_2)=1$. The fourth term of \eqref{com1} is given by
	\begin{eqnarray}
		\gamma_2(a_1, a_2,b) &= & \nu_b(f(\theta_2(b), a_1a_2)\theta_1(a_1 a_2))(\theta_1(\beta_b(a_1 a_2)))^{-1} \nu_{b}(\theta_1(a_2)(\theta_1(a_1 a_2))^{-1}\mu_{a_2}(\theta_1(a_1))) \notag\\
		&&  (\mu_{ \beta_{b}(a_2)}(\nu_b(f(\theta_2(b), a_1)\theta_1(a_1))(\theta_1(\beta_b(a_1)))^{-1}))^{-1} (\nu_b(f(\theta_2(b), a_2)\theta_1(a_2)) (\theta_1(\beta_b(a_2)))^{-1})^{-1}\notag\\
		&&(\theta_1(\beta_{b}(a_2)))^{-1}\theta_1(\beta_{b}(a_1)\beta_{b}(a_2))(\mu_{\beta_{b}(a_2)}(\theta_1(\beta_{b}(a_1))))^{-1}.\notag
	\end{eqnarray}
	Using that $f(\theta_2(b), a_1 a_2) =  \mu_{ a_2}(f(\theta_2(b),a_1))f(\theta_2(b), a_2)$ and $\nu_b(\mu_{ a_2}(k))=\mu_{ \beta_{b}(a_2)}(\nu_b(k))$ for all $k \in K$, we have $\gamma_2(a_1, a_2,b)=1$. The fifth term is given by 
	\begin{eqnarray*}
		\gamma_3(a_1, a_2) &= & \chi(a_1 \circ_T a_2) ~S(\nu^{-1}_{T(a_1 \circ a_2)}\big(\rho(a_2, T(a_1))\tau_1(a_1,\beta_{T(a_1)}(a_2)) \nu_{T(a_1)}(f(\chi(a_1), a_2)))\big)\notag\\
		&& (\sigma_{T(a_2)}( \chi(a_1)))^{-1} ~(\chi(a_2))^{-1} ~(\tau_2(T(a_1), T(a_2) ))^{-1} \notag\\
		&= &  S(\nu^{-1}_{T(a_1 \circ_T a_2)}(\theta_1(a_1 \circ_T a_2))) (\theta_2(T(a_1 \circ_T a_2)))^{-1} S(\nu^{-1}_{T(a_1 \circ_T a_2)}\big( \nu_{T(a_1)}(f(\theta_2(T(a_1)), a_2)\notag\\
		&& \theta_1(a_2))(\theta_1(\beta_{T(a_1)}(a_2)))^{-1} \theta_1(\beta_{T(a_1)}(a_2))(\theta_1(a_1 \beta_{T(a_1)}(a_2)))^{-1}\mu_{\beta_{T(a_1)}(a_2)}(\theta_1(a_1))\notag\\
		&&  \nu_{T(a_1)}(f(S(\nu^{-1}_{T(a_1)}(\theta_1(a_1)))(\theta_2(T(a_1)))^{-1}, a_2 ))) \big)  (\sigma_{T(a_2)}(S(\nu^{-1}_{T(a_1)}(\theta_1(a_1)))(\theta_2(T(a_1)))^{-1}))^{-1}\notag\\   
		&&	\big(S(\nu^{-1}_{T(a_2)}(\theta_1(a_2)))(\theta_2(T(a_2)))^{-1}\big)^{-1} \big( \theta_2(T(a_2))(\theta_2(T(a_1) T(a_2)))^{-1}\sigma_{T(a_2)}(\theta_2(T(a_1))) \big)^{-1}.
	\end{eqnarray*}
	Using that $\nu^{-1}_{T(a_1)}(\mu_{\beta_{T(a_1)}(a_2)}(\theta_1(a_1)))= \mu_{a_2}(\nu^{-1}_{T(a_1)}(\theta_1(a_1)))$, the only part that survive is  
	\begin{align*}
		S\big(\nu^{-1}_{T(a_2)} \big(f(S(\nu^{-1}_{T(a_1)}(\theta_1(a_1))), a_2) \big)\big)\, S\big(\nu^{-1}_{T(a_2)} \big(\mu_{a_2}(\nu^{-1}_{T(a_1)}(\theta_1(a_1))) \big)\big) \,\big(\sigma_{T(a_2)} \big(S(\nu^{-1}_{T(a_1)}(\theta_1(a_1))) \big)\big)^{-1},
	\end{align*}
	which is 1 from \eqref{Rallcondn}  by putting $\nu^{-1}_{T(a_1)}(\theta_1(a_1))=k$. This completes the proof of the lemma.
\end{proof}

\end{document}
